% DhatuChakra — Technical Documentation
% Compile: pdflatex (Overleaf: just upload this file and click Recompile)
\documentclass[11pt,a4paper]{article}

\usepackage[margin=2.4cm]{geometry}
\usepackage[T1]{fontenc}
\usepackage[utf8]{inputenc}
\usepackage{lmodern}
\usepackage{textcomp}
\usepackage{amsmath}
\usepackage{booktabs}
\usepackage{array}
\usepackage{xcolor}
\usepackage{enumitem}
\usepackage{titlesec}
\usepackage[hidelinks,colorlinks=true,linkcolor=navy,urlcolor=saffron]{hyperref}
\usepackage{fancyhdr}

% Brand palette
\definecolor{saffron}{HTML}{B84E0A}
\definecolor{navy}{HTML}{24466E}
\definecolor{green}{HTML}{17803D}
\definecolor{inksoft}{HTML}{5A6B7B}

\titleformat{\section}{\Large\bfseries\color{navy}}{\thesection}{0.8em}{}
\titleformat{\subsection}{\large\bfseries\color{navy}}{\thesubsection}{0.8em}{}

\pagestyle{fancy}
\fancyhf{}
\fancyhead[L]{\small\color{inksoft}DhatuChakra --- Technical Documentation}
\fancyhead[R]{\small\color{inksoft}Smart India Hackathon $\cdot$ Problem \#48}
\fancyfoot[C]{\small\thepage}

\setlength{\parskip}{4pt}
\setlength{\parindent}{0pt}

\begin{document}

% ---------------- Title ----------------
\begin{center}
  {\Huge\bfseries\color{saffron} DhatuChakra}\\[2pt]
  {\large\color{inksoft} (Dhatu = metal $\cdot$ Chakra = the wheel)}\\[8pt]
  {\LARGE\bfseries AI-Driven Life Cycle Assessment \& Circularity\\Engine for Indian Metals}\\[10pt]
  {\normalsize Technical Documentation --- Working Prototype v1}\\[2pt]
  {\normalsize\color{inksoft} Smart India Hackathon (Internal Round) $\cdot$ Problem Statement \#48 $\cdot$ Ministry of Mines}\\[2pt]
  {\normalsize\color{inksoft} Repository: \url{https://github.com/grv-io/LCA-TOOL}}
\end{center}

\vspace{4pt}
\hrule
\vspace{10pt}

\tableofcontents

% ================================================================
\section{Problem and Motivation}

Life Cycle Assessment (LCA) is the standard method for quantifying the environmental
footprint of metal production, and circularity indicators are becoming central to both
national policy (National Critical Mineral Mission, circular-economy frameworks) and
trade (the EU Carbon Border Adjustment Mechanism). In practice, however, LCA is out of
reach for most Indian plants:

\begin{itemize}[leftmargin=1.4em, itemsep=1pt]
  \item a conventional study needs \textbf{25+ process parameters} most plants do not track;
  \item commercial inventory databases (e.g.\ ecoinvent) cost $\approx$ \texteuro 4{,}000 per year;
  \item consultant-led studies take weeks and are not comparable across plants;
  \item global tools ignore Indian conditions (grid intensity, state-level variation, domestic data sources).
\end{itemize}

\textbf{DhatuChakra} answers this with a tool that needs only 5--8 inputs (or one plain
sentence), fills the remaining parameters with a machine-learning imputer that reports its
own confidence, and returns a full impact and circularity profile with quantified
uncertainty --- every number traceable to a cited public source.

% ================================================================
\section{System Overview}

\subsection{What the prototype does}
\begin{enumerate}[leftmargin=1.6em, itemsep=1pt]
  \item \textbf{Scenario input} --- metal (aluminium, steel, copper), production route,
        region or Indian state grid, recycled content, end-of-life recovery, renewable
        electricity share; optionally a plain-English description parsed by keyword matching.
  \item \textbf{Parameter estimation} --- a $k$-nearest-neighbours model predicts electricity
        and thermal-energy intensity with a real confidence score; every estimated field is
        visibly flagged and user-overridable.
  \item \textbf{Impact computation} --- GWP (kg CO$_2$e), cumulative energy (GJ), water
        (m$^3$) and acidification (kg SO$_2$e) per tonne of metal, cradle-to-gate.
  \item \textbf{Uncertainty} --- a 1{,}000-run Monte Carlo simulation produces p05--p95
        ranges and a histogram.
  \item \textbf{Circularity} --- the Material Circularity Indicator (Ellen MacArthur
        Foundation methodology), rendered as a 24-spoke chakra gauge.
  \item \textbf{Decision support} --- linear-vs-circular comparison with live sliders,
        benchmark strip (your plant vs India average vs world best practice), CBAM
        export-cost calculator, rule-based recommendations that cite the user's own numbers,
        and an ISO-14044-style printable report.
\end{enumerate}

\subsection{Architecture}

The prototype is a fully client-side web application (HTML/CSS/JavaScript, no backend, no
external libraries; runs offline). The calculation engine, ML model and all data tables are
transparent, inspectable files:

\begin{center}
\begin{tabular}{@{}ll@{}}
\toprule
\textbf{Module} & \textbf{Responsibility} \\
\midrule
\texttt{js/data.js}    & All constants: routes, grids, benchmarks, CBAM, presets (single source of truth) \\
\texttt{js/calc.js}    & Engine: impact formulas, MCI, Monte Carlo, recommendations, text parser \\
\texttt{js/ml\_data.js} & 500 physics-generated training rows for the $k$-NN imputer \\
\texttt{js/ui.js}      & Rendering: chakra gauge, Sankey, stage bar, histogram, benchmarks, Hindi toggle \\
5 HTML pages           & Dashboard, Assess, Results, Compare, Report \\
\bottomrule
\end{tabular}
\end{center}

Flow: \emph{user input} $\rightarrow$ \emph{parser} $\rightarrow$ \emph{$k$-NN imputer
(+confidence)} $\rightarrow$ \emph{LCA engine $\times$ factor library} $\rightarrow$
\emph{Monte Carlo} $\rightarrow$ \emph{results, MCI, benchmarks, CBAM, recommendations,
report}.

The full production build (designed, archived in the repository) replaces the factor engine
with Brightway2.5, the $k$-NN with quantile-gradient-boosting imputation, and the keyword
parser with an LLM parser --- behind a FastAPI backend with PostgreSQL.

% ================================================================
\section{Calculation Methodology}

\subsection{Functional unit and boundary}
The functional unit is \textbf{one tonne of metal at plant gate}; the system boundary is
\textbf{cradle-to-gate}. The model is attributional.

\subsection{Impact model}
For each metal a linear (primary) and circular (secondary) route is parameterised by
electricity intensity $E$ (kWh/t), non-electric base emissions $B$ (kg CO$_2$e/t) and base
energy $G$ (GJ/t). With recycled content $r \in [0,1]$, renewable electricity share
$s \in [0,1]$ and grid intensity $g_0$ (kg CO$_2$e/kWh):

\begin{align}
g_{\mathrm{eff}} &= g_0\,(1-s) + 0.04\,s\\
\mathrm{GWP} &= (1-r)\,\big(E_{\mathrm{lin}}\, g_{\mathrm{eff}} + B_{\mathrm{lin}}\big)
              + r\,\big(E_{\mathrm{circ}}\, g_{\mathrm{eff}} + B_{\mathrm{circ}}\big)\\
\mathrm{Energy} &= (1-r)\,\big(0.0103\,E_{\mathrm{lin}} + G_{\mathrm{lin}}\big)
                 + r\,\big(0.0103\,E_{\mathrm{circ}} + G_{\mathrm{circ}}\big)\\
\mathrm{Water} &= 0.004 \cdot \mathrm{GWP} \qquad
\mathrm{Acidification} = 0.005 \cdot \mathrm{GWP}
\end{align}

where $0.04$ is the life-cycle intensity of the renewable mix and $0.0103$ GJ/kWh the
primary-energy factor including generation losses. Route parameters:

\begin{center}
\begin{tabular}{@{}llrrr@{}}
\toprule
\textbf{Metal} & \textbf{Route} & $E$ (kWh/t) & $B$ (kg CO$_2$e/t) & $G$ (GJ/t) \\
\midrule
Aluminium & Primary (Hall--H\'eroult) & 14{,}500 & 6{,}300 & 50 \\
Aluminium & Recycled (remelt)         & 800      & 350     & 8  \\
Steel     & BF--BOF                   & 120      & 2{,}100 & 21 \\
Steel     & EAF (scrap)               & 450      & 380     & 2  \\
Copper    & Primary (pyromet.)        & 3{,}300  & 1{,}900 & 28 \\
Copper    & Recycled (secondary)      & 900      & 500     & 6  \\
\bottomrule
\end{tabular}
\end{center}

$B$ for primary aluminium contains anode CO$_2$, PFC emissions, alumina refining thermal
energy and mining/calcination/transport; for BF--BOF steel it contains coke, blast furnace
and converter emissions (hence the small electricity term).

\subsection{Grid intensity}
National factors: India 0.71 (CEA CO$_2$ Baseline Database v19), world average 0.48,
Europe 0.28 kg CO$_2$e/kWh. Additionally, ten \emph{indicative} Indian state factors
(derived from CEA v19 plus state generation mix) drive the interactive state map, from
Himachal 0.18 (hydro) to Chhattisgarh 0.92 (coal). The same smelter moves from
$\approx$8.9 to $\approx$19.1 t CO$_2$e/t between Himachal and Odisha --- plant location is
a climate decision.

\subsection{Material Circularity Indicator}
Following the Ellen MacArthur Foundation methodology (v3), with recycled feedstock fraction
$F_r = r$, end-of-life collection rate $C_r$, recycling process efficiency $E_f = 0.9$ and
utility factor $X = 1$ (mass normalised to $M=1$):

\begin{align}
V &= 1 - F_r, \qquad W_0 = 1 - C_r, \qquad W_c = C_r\,(1-E_f), \qquad W_f = \frac{F_r\,(1-E_f)}{E_f}\\
W &= W_0 + \tfrac{1}{2}(W_f + W_c), \qquad
\mathrm{LFI} = \frac{V + W}{2 + \tfrac{1}{2}(W_f - W_c)}, \qquad
\mathrm{MCI} = \max\!\big(0,\; 1 - 0.9\,\mathrm{LFI}\big)
\end{align}

The UI renders MCI as an Ashoka-chakra gauge: $\mathrm{round}(24 \cdot \mathrm{MCI})$ of 24
spokes fill green.

\subsection{Monte Carlo uncertainty}
Each displayed range is the p05--p95 interval of 1{,}000 simulation runs. Per run,
independent lognormal noise (geometric standard deviation 1.1) multiplies the grid factor,
electricity intensity and base emissions:
$x' = x \cdot \exp\!\big(\ln(1.1)\, Z\big)$, $Z \sim \mathcal{N}(0,1)$ via Box--Muller.
Water and acidification ranges derive from the GWP samples through their fixed ratios. The
histogram is drawn from the same samples and is recomputed on every load.

\subsection{$k$-NN parameter imputer}
The ``estimated'' fields on the assessment page are predictions of a $k$-nearest-neighbours
model, not hardcoded values.

\begin{itemize}[leftmargin=1.4em, itemsep=1pt]
  \item \textbf{Training data:} 500 rows generated from the engine's own physics with
        $\pm 8\%$ noise (script \texttt{scripts/make\_training\_data.js}); regenerable at
        any time.
  \item \textbf{Features (normalised):} metal indicator, recycled content $r$, renewable
        share $s$, grid factor.
  \item \textbf{Prediction:} mean of the $k=7$ nearest neighbours (Euclidean distance) for
        electricity and thermal intensity.
  \item \textbf{Confidence:} $1 - \sigma/\mu$ over the neighbours, clamped to $[0.50, 0.99]$
        --- it genuinely drops in uncertain regions (e.g.\ 94\% at pure primary aluminium
        vs 78\% at a 50\% scrap blend).
\end{itemize}

The imputer is a suggestion layer: the engine itself always uses library physics values
unless the user accepts or edits an estimate (auditable separation of AI and physics).
Copper currently falls back to library defaults (model trained on Al/steel; retraining with
copper is scheduled).

\subsection{CBAM export-cost model}
For EU-bound exports, annual CBAM exposure is estimated as
\begin{equation}
\mathrm{Cost} = \max\!\big(0,\; \mathrm{GWP}_t - b_m\big) \cdot \phi_y \cdot P_{\mathrm{ETS}}
\cdot X_{\mathrm{EUR \to INR}} \cdot T
\end{equation}
with indicative EU benchmark $b_m$ (Al 6.0, steel 1.8 t CO$_2$e/t), phase-in factor
$\phi_y$ per Regulation (EU) 2023/956 (2026: 2.5\%, 2027: 5\%, 2028: 10\%, 2029: 22.5\%,
2030: 48.5\%, 2034: 100\%), adjustable ETS price $P_{\mathrm{ETS}}$ (default \texteuro 75/t)
and export tonnage $T$. Example: a primary-aluminium smelter (16.6 t CO$_2$e/t) exporting
10{,}000 t/yr faces $\approx$ Rs.~1.8 crore in 2026, rising to $\approx$ Rs.~35 crore
by 2030. Copper is outside the current CBAM product scope; the tool states this instead of
showing a cost.

% ================================================================
\section{Validation}

Engine outputs are pinned to an expected-output table; any deviation is treated as a bug.
All six presets sit inside published ranges (IAI, worldsteel, EF 3.1 summaries):

\begin{center}
\begin{tabular}{@{}lrrrrrr@{}}
\toprule
\textbf{Preset (India grid)} & $r$ & $C_r$ & \textbf{GWP} & \textbf{Energy} & \textbf{Water} & \textbf{MCI} \\
 & & & (t CO$_2$e/t) & (GJ/t) & (m$^3$/t) & \\
\midrule
Aluminium --- Primary  & 0.00 & 0.15 & 16.6 & 199  & 66   & 0.16 \\
Aluminium --- Recycled & 1.00 & 0.70 & 0.92 & 16   & 3.7  & 0.83 \\
Steel --- BF--BOF      & 0.00 & 0.25 & 2.19 & 22   & 8.7  & 0.20 \\
Steel --- EAF          & 1.00 & 0.85 & 0.70 & 6.6  & 2.8  & 0.89 \\
Copper --- Primary     & 0.00 & 0.30 & 4.24 & 62   & 17.0 & 0.22 \\
Copper --- Recycled    & 1.00 & 0.80 & 1.14 & 15.3 & 4.6  & 0.87 \\
\bottomrule
\end{tabular}
\end{center}

Hand check (Al primary): $14{,}500 \times 0.71 + 6{,}300 = 16{,}595$ kg $\approx 16.6$ t
CO$_2$e/t --- inside the IAI-consistent 16--20 range for coal-heavy grids. The recycled
route is $94\%$ lower, matching the widely cited $\approx 95\%$ saving. Reference ranges:
recycled Al 0.5--1.5, BF--BOF 1.9--2.4, EAF 0.4--1.0 t CO$_2$e/t.

Benchmarks shown in the ``Where you stand'' card (indicative): Al primary --- India average
17.8, world best practice (hydro smelters) 5.5; steel BF--BOF --- India 2.55, best 1.8;
copper primary --- India 4.6, best 3.0 t CO$_2$e/t.

% ================================================================
\section{Data Sources}

\begin{center}
\begin{tabular}{@{}lll@{}}
\toprule
\textbf{Factor} & \textbf{Source} & \textbf{Year} \\
\midrule
India grid intensity & CEA CO$_2$ Baseline Database v19 & 2023 \\
State grid intensities (indicative) & CEA v19 + state generation mix & 2023 \\
World/EU grid intensity & IEA emission factors & 2023 \\
Al smelting electricity, PFC, anode & International Aluminium Institute LCI & 2022 \\
Alumina refining and mining & EF 3.1 / ELCD & 2022 \\
Al remelting route & IDEMAT 2024, TU Delft & 2024 \\
BF--BOF and EAF inventories & worldsteel LCI methodology & 2023 \\
Copper routes (indicative) & EF 3.1 / IDEMAT 2024 & 2024 \\
Transport factors & EF 3.1 & 2022 \\
MCI methodology & Ellen MacArthur Foundation, MCI v3 & 2019 \\
CBAM schedule and scope & Regulation (EU) 2023/956 & 2023 \\
\bottomrule
\end{tabular}
\end{center}

Every factor is public and free --- the tool has \textbf{zero database licensing cost}, and
each on-screen figure links back to its source (provenance table in the app and in the
generated report).

% ================================================================
\section{Product Features (Prototype v1)}

\begin{itemize}[leftmargin=1.4em, itemsep=1pt]
  \item \textbf{Dashboard} with six preset scenarios and a saved-scenario library
        (save/load/delete, persists locally).
  \item \textbf{Assess} --- form + plain-English entry; interactive India state-grid map;
        $k$-NN estimated fields with live confidence badges; every estimate overridable.
  \item \textbf{Results} --- four KPI cards with Monte-Carlo p05--p95 ranges and histogram;
        stage-contribution bar; material-flow Sankey with the recycling return loop; chakra
        MCI gauge; benchmark strip; CBAM calculator; provenance table.
  \item \textbf{Compare} --- locked linear baseline vs live circular scenario; three
        circularity sliders; savings delta strip; recommendations that cite the current
        numbers.
  \item \textbf{Report} --- ISO-14044-style document (goal \& scope, inventory, results with
        uncertainty, circularity, recommendations, sources, limitations); print-to-PDF.
  \item \textbf{Hindi/English toggle} for the static UI ($\approx$54 strings; full
        localisation planned).
  \item Fully offline-capable: fonts, map and all data ship with the page.
\end{itemize}

% ================================================================
\section{Limitations and Honest Caveats}

\begin{enumerate}[leftmargin=1.6em, itemsep=1pt]
  \item The engine is a factor model with fixed water/acidification ratios --- adequate for
        route comparison, not for site permitting. The production design swaps in a full
        matrix LCA (Brightway2.5).
  \item State grid factors and copper/benchmark values are \emph{indicative}, clearly
        labelled as such, pending plant-level data.
  \item The $k$-NN imputer is trained on physics-generated (not measured) data; it
        demonstrates the imputation architecture and honest confidence reporting rather
        than field accuracy. Copper is not yet in the training set.
  \item The plain-English parser is keyword-based in the prototype; the production build
        uses an LLM parser with validation.
  \item Cradle-to-gate only; use-phase and end-of-life stages are parameterised in the MCI
        but not in the impact totals.
\end{enumerate}

% ================================================================
\section{Roadmap (Full Build)}

FastAPI + PostgreSQL backend; Brightway2.5 matrix engine; 100+ row cited factor library
with regional fallback chains; LightGBM quantile imputer (p10/p50/p90) and surrogate model
(target $R^2 > 0.95$, $<$20 ms); LLM scenario parser and grounded recommendations;
Monte-Carlo with per-factor pedigree uncertainty; multi-user scenario storage; full Hindi
localisation; validation notebook against 8--10 published studies. The complete engineering
plan, interface contracts and per-module specifications are archived in the repository
(\texttt{docs/archive/}).

% ================================================================
\section*{References}
\begin{enumerate}[leftmargin=1.6em, itemsep=0pt]
  \item ISO 14040/14044 --- Environmental management, LCA: principles and requirements. \url{https://www.iso.org}
  \item Ellen MacArthur Foundation, \emph{Material Circularity Indicator methodology} (v3). \url{https://www.ellenmacarthurfoundation.org}
  \item International Aluminium Institute --- LCI data and GHG pathways. \url{https://international-aluminium.org}
  \item worldsteel --- Life Cycle Inventory methodology report. \url{https://worldsteel.org}
  \item Central Electricity Authority --- CO$_2$ Baseline Database v19. \url{https://cea.nic.in}
  \item EU Environmental Footprint EF 3.1 / ELCD. \url{https://eplca.jrc.ec.europa.eu}
  \item IDEMAT 2024, TU Delft. \url{https://www.idematapp.com}
  \item Indian Bureau of Mines --- Indian Minerals Yearbook. \url{https://ibm.gov.in}
  \item Regulation (EU) 2023/956 (CBAM). \url{https://taxation-customs.ec.europa.eu}
  \item NITI Aayog --- circular-economy strategy papers. \url{https://www.niti.gov.in}
\end{enumerate}

\vspace{10pt}
\hrule
\vspace{6pt}
\begin{center}
\small\color{inksoft}
DhatuChakra $\cdot$ every tonne, accounted for $\cdot$ built for Smart India Hackathon, Problem \#48\\
Working prototype: all formulas, data and the ML model in this document are implemented and verifiable in the repository.
\end{center}

\end{document}
